% !TEX encoding = UTF-8 Unicode
% !TEX root = thesis-main.tex

\chapterr{\textlatin{MeL}}
Στο προηγούμενο κεφάλαιο αναλύσαμε το σκοπό και αναφέραμε επιγραμματικά τις υπηρεσίες του συστήματος. Το \textlatin{MeL} έχει ως κεντρική ιδέα την εκπαίδευση των χρηστών καθώς και την επικοινωνία μεταξύ τους. Παρέχει τη δυνατότητα σε χρήστες με δικαιώματα διαχειριστή τη δημιουργία και την καταχώρηση εκπαιδευτικών άρθρων και τεστ στα οποία έχουν πρόσβαση οι χρήστες-μέλη για την ανάγνωση και την επίλυσή τους αντίστοιχα.
Αυτό το κεφάλαιο έχει ως σκοπό την αναφορά και την παρουσίαση των προγραμμάτων και των γλωσσών που χρησιμοποιήθηκαν, την ανάλυση της Βάσης Δεδομένων και τέλος την αναφορά των απαιτήσεων για την εγκατάσταση του \textlatin{MeL}.

\section{Αρχιτεκτονική συστήματος}
Για την ανάπτυξη του συστήματος χρησιμοποιήσαμε το πλαίσιο \textlatin{Windows/apache/postgreSQL/PHP (WAPP)} σαν υποδομή της βάσης δεδομένων ενώ στα υπόλοιπα τμήματα χρησιμοποιήθηκαν τα εργαλεία \textlatin{HTML, Javascript, PHP, jQuery, CSS.} Οι λόγοι επιλογής αυτών για την υλοποίηση του \textlatin{MeL} είναι γιατί και τα τρία προϊόντα είναι ανοιχτού κώδικα, συνεργάζονται μεταξύ τους με απόλυτη επιτυχία, υπάρχουν πάρα πολλές συναρτήσεις της \textlatin{PHP} για τη \textlatin{PSQL} και τέλος η συνδυασμένη τους χρήση επιτρέπει τη δημιουργία δυναμικών διαδικτυακών εφαρμογών με υψηλά κριτήρια αξιοπιστίας, απόδοσης, ταχύτητας και ασφάλειας. 

Ο τρόπος με τον οποίο τα τρία βασικότερα προγράμματα της εφαρμογής συνδέονται και επικοινωνούν μεταξύ τους περιγράφεται με απλά λόγια ως εξής: 
\begin{itemize}
\item
 Αρχικά ο \textlatin{browser} του χρήστη στέλνει ένα ερώτημα \textlatin{HTTP (HTTP Request)} προς τον \textlatin{Apache server} σχετικά με τη σελίδα που προσπαθεί να ανοίξει. 
\item
 Ο \textlatin{Apache} εντοπίζει το αρχείο με τον κατάλληλο κώδικα και καλεί τον μεταγλωττιστή της \textlatin{PHP} για να εκτελέσει τον κώδικα. 
\item
 Η \textlatin{PHP} μέσω ενσωματωμένων συναρτήσεων και λειτουργιών που διαθέτει στέλνει ερωτήματα \textlatin{(queries)} προς τη βάση δεδομένων \textlatin{psql.} 
\item
 Η \textlatin{psql} επεξεργάζεται τα ερωτήματα και στέλνει πίσω στην \textlatin{PHP} τα αποτελέσματα. 
\item
 Η \textlatin{PHP} επιστρέφει πίσω στον \textlatin{Apache server} τα αποτελέσματα της βάσης δεδομένων και του κώδικα σε μορφή \textlatin{HTML}. 
\item
 Ο \textlatin{Apache server} στέλνει την \textlatin{HTML} στον \textlatin{browser} όπου προβάλλονται στο χρήστη. 
\end{itemize}

\subsection{\textlatin{Apache Server}}
Ο \textlatin{Apache HTTP \cite{apache,apache_wiki}} ...

\subsection{\textlatin{PostgreSQL}}
Η \textlatin{PostgreSQL \cite{postgres}} ...

\subsection{\textlatin{HTML}}
Η γλώσσα \textlatin{HTML \cite{html} (Hypertext Markup Language)} ...

Δημιουργήθηκε με σκοπό: 
\begin{itemize}
\item
 Να μειωθεί η εξάρτηση από εξωτερικές τεχνολογίες όπως το \textlatin{Adobe Flash.} 
\item
 Τα χαρακτηριστικά της να είναι βασισμένα στις τεχνολογίες \textlatin{HTML, CSS, DOM} και \textlatin{JavaScript.} 
\item
 Να είναι η \textlatin{HTML5} ανεξάρτητη από την συσκευή στην οποία εκτελείται.
\item
 Καλύτερο χειρισμό λαθών. 
\item
 Να χρησιμοποιείται περισσότερη σήμανση παρά \textlatin{scripts.} 
\end{itemize}


\subsection{\textlatin{PHP}}
Η \textlatin{PHP (Hypertext Preprocessor)} \cite{www11} ...

\subsection{\textlatin{CSS}}
Η γλώσσα \textlatin{CSS} \cite{www11} ...

\subsection{\textlatin{jQuery}}
Η \textlatin{jQuery \cite{jquery}} ...


\section{Βάση Δεδομένων}
Στο κεφάλαιο αυτό γίνεται αναφορά στη Βάση Δεδομένων πάνω στην οποία στηρίζεται το \textlatin{MeL}.  Αρχικά, παρουσιάζεται το σχήμα της Βάσης Δεδομένων, έπειτα γίνεται αναφορά και περιγραφή των πινάκων από τους οποίους αποτελείται και τέλος παρουσιάζονται οι εντολές που εκτελούνται για τη δημιουργία της.

\subsection{Σχήμα της Βάσης Δεδομένων}
Στην παράγραφο αυτή θα παρουσιάσουμε το σχήμα της Βάσης Δεδομένων. Από τις απαιτήσεις, τους περιορισμούς αλλά και τα χαρακτηριστικά της Βάσης Δεδομένων δημιουργείται το παρακάτω σχήμα. Στο Σχήμα \ref{fig:schemaDB}, παρατηρούμε τους πίνακες οι οποίοι απαρτίζουν την Βάση Δεδομένων του \textlatin{MeL.} Στους πίνακες αυτούς αποθηκεύονται όλα τα δεδομένα της πλατφόρμας όπως, τα στοιχεία των χρηστών, μηνύματα, άρθρα, τεστ, λύσεις των τεστ κ.α. 


\begin{figure}[t]
% \centerline{\includegraphics[scale=0.5]{figures/ch2-db_schema_3}}
\caption{Σχήμα Βάσης Δεδομένων \textlatin{MeL.}}
\label{fig:schemaDB}
\end{figure}

\subsection{Πίνακες Βάσης Δεδομένων}
Στην παράγραφο αυτή θα παρουσιάσουμε τους πίνακες από τους οποίους αποτελείται η Βάση Δεδομένων της πλατφόρμας. Θα αναφέρουμε ξεχωριστά για τον καθένα τα δεδομένα που αποθηκεύει και θα παραθέσουμε την εντολή που πρέπει να εκτελεστεί για την δημιουργία του.

\begin{table}[t]
\centering
\begin{tabular}{|m{2.5cm}|m{2cm}|m{9cm}|}
\hline
\rowcolor{gray} 
ΠΕΔΙΟ & ΤΥΠΟΣ & ΠΕΡΙΓΡΑΦΗ \\
\hline\hline
\textlatin{user\_id} & \textlatin{integer} & Το \textlatin{id} του χρήστη. \\ 
\textlatin{username} & \textlatin{varchar} & Το συνθηματικό όνομα του χρήστη. \\
\textlatin{password} & \textlatin{varchar} & Ο κωδικός πρόσβασης του χρήστη.  \\ 
\textlatin{first\_name} & \textlatin{varchar} & Το πραγματικό όνομα του χρήστη. \\ 
\textlatin{last\_name} & \textlatin{varchar} & Το επώνυμο του χρήστη. \\ 
\textlatin{email} & \textlatin{varchar} & Το \textlatin{email} του χρήστη. \\ 
\textlatin{email\_code} & \textlatin{varchar} & Κωδικός που δημιουργείται τυχαία και χρησιμοποιείται για την επιβεβαίωση της εγγραφής. \\
\textlatin{active} & \textlatin{integer} & Μεταβλητή για το αν είναι ενεργοποιημένος ο λογαριασμός του χρήστη ή όχι. Παίρνει αρχική τιμή 0 και μόλις ενεργοποιηθεί ο λογαριασμός του παίρνει την τιμή 1. \\ 
\textlatin{rights} & \textlatin{varchar} & Μεταβλητή για τα δικαιώματα που έχει ο χρήστης στην πλατφόρμα. Παίρνει αρχική τιμή \textlatin{user} και μπορεί να πάρει τις τιμές \textlatin{admin} ή \textlatin{superadmin.}\\
\textlatin{notifications} & \textlatin{integer} & Μεταβλητή που μετράει το πλήθος των ειδοποιήσεων του χρήστη. Παίρνει αρχική τιμή 0. \\ 
\hline
\end{tabular}
\caption{Πίνακας Βάσης Δεδομένων \textlatin{users}}
\label{table:users}
\end{table}



Ο Πίνακας \ref{table:users} παρουσιάζει τον πίνακα της Βάσης Δεδομένων \textlatin{users} που αποθηκεύει τα βασικά στοιχεία των χρηστών για τη δημιουργία λογαριασμού στην πλατφόρμα. Για να δημιουργήσουμε τον πίνακα \textlatin{users} αρχικά δημιουργούμε μια ακολουθία την οποία θα χρησιμοποιήσουμε για παίρνει μοναδικές τιμές το \textlatin{id} του χρήστη που είναι και κλειδί του πίνακα. Η εντολή για την ακολουθία αυτή είναι η εξής:
{\latintext
\begin{lstlisting}[language=SQL]
CREATE SEQUENCE user_id_seq
    START WITH 1
    INCREMENT BY 1
    NO MINVALUE
    NO MAXVALUE
    CACHE 1;
\end{lstlisting}
}
και στη συνέχεια εκτελούμε την εντολή για τον πίνακα η οποία είναι:
{\latintext
\begin{lstlisting}[language=SQL]
CREATE TABLE users (
    user_id integer DEFAULT
    nextval('user_id_seq'::regclass) NOT NULL,
    username character varying(32),
    password character varying(32),
    first\_name character varying(32),
    last\_name character varying(32),
    email character varying(1024),
    email\_code character varying(32),
    active integer DEFAULT 0,
    rights character varying(32) DEFAULT 'user'::character varying,
    notifications integer DEFAULT 0,
    PRIMARY KEY (user\_id)
);
\end{lstlisting}
}


%
%
%
\begin{algorithm}[t]\latintext
\Algorithm{\textsc{Frequent\-Sub\-graph\-Mining}}
\scriptsize
\BlankLine

\KwIn{A graph $G$ and the frequency threshold $\tau$} 
\KwOut{All subgraphs $S$ of $G$ such that $s_G(S)\geq \tau$} 
\BlankLine

$\mathit{result} \leftarrow\emptyset$\\
Let $\mathit{fEdges}$ be the set of all frequent edges of $G$ %(i.e., with support $\geq\tau$)\\
%Sort $S$ in DFS lexicographical order

\ForEach{$e \in \mathit{fEdges}$}
{
	$\mathit{result} \leftarrow \mathit{result} \cup \textsc{SubgraphExtension}(e, G, \tau, \mathit{fEdges})$  \\
	Remove $e$ from $G$ and $\mathit{fEdges}$
}
\Return $\mathit{result}$
\end{algorithm}
%
%
%
\begin{algorithm}[t]\latintext
\Algorithm{\textsc{SubgraphExtension}}
\BlankLine

\scriptsize
\KwIn{A subgraph $S$ of a graph data $G$, the frequency threshold $\tau$ and the set of frequent edges $\mathit{fEdges}$ of $G$} 
\KwOut{All frequent subgraphs of $G$ that extend $S$} 
\BlankLine

$\mathit{result} \leftarrow S$, $\mathit{candidateSet} \leftarrow \emptyset$

\ForEach{edge $e$ in $\mathit{fEdges}$ and node $u$ of $S$} {\label{lin:SE-for1-s}
		\If{$e$ can be used to extend $u$} {
			Let $ext$ be the extension of $S$ with $e$\\
			\lIf{$ext$ is not already generated}{
				$\mathit{candidateSet}\leftarrow\mathit{candidateSet}\cup ext$
				\label{lin:SE-for1-e}
			}
		}
	}


\ForEach{$c \in \mathit{candidateSet}$}
{ \label{lin:SE-for2-s}
	\If{$s_G(c)\geq\tau$ \label{lin:freq-eval} }
	{
		$\mathit{result}\leftarrow\mathit{result}\cup \textsc{SubgraphExtension}(c, G, \tau, \mathit{fEdges})$
		\label{lin:SE-for2-e}
	}
}
\Return $\mathit{result}$
\end{algorithm}

%  
%
\begin{algorithm}[t]\latintext
\Algorithm{\textsc{IsFre\-que\-nt\-Csp}}
\scriptsize
\BlankLine

\KwIn{Graphs $S$ and $G$ and the frequency threshold $\tau$} 
\KwOut{$\mathit{true}$ if $S$ is a frequent subgraph of $G$, $\mathit{false}$ otherwise} 
\BlankLine

Consider the subgraph $S$ to graph $G$ CSP\\
Apply node and arc consistency \label{lin:csp-arc-consistency}\\
\lIf{the size of any domain is less than $\tau$}{\Return $\mathit{false}$}\label{lin:dom-less-csp}
\ForEach{solution $\mathcal{S}ol$ of the $S$ to graph $G$ CSP}{
	Mark all nodes of $\mathcal{S}ol$ in the corresponding domains \label{lin:isFreq-csp}\\
	\lIf{all domains have at least $\tau$ nodes marked} 
		{\Return $\mathit{true}$}\label{lin:isFreq-domdon-csp}
}
\Return $\mathit{false}$\ \ \  \tcp{Domain is exhausted}
\end{algorithm}


%
%ISFREQUENT
%
\begin{algorithm}[t]\latintext
\Algorithm{\textsc{IsFre\-qu\-ent}}
\scriptsize
%\small
\BlankLine

%\KwIn{A subgraph $S$ of a graph $G$ and the frequency threshold $\tau$} 
%\KwOut{$\mathit{true}$ if $S$ is frequent in $G$, $\mathit{false}$ otherwise} 
\KwIn{Graphs $S$ and $G$ and the frequency threshold $\tau$} 
\KwOut{$\mathit{true}$ if $S$ is a frequent subgraph of $G$, $\mathit{false}$ otherwise} 
\BlankLine

Consider the subgraph $S$ to graph $G$ CSP and apply node and arc consistency\\ 


\BlankLine
\Comment{Push-down pruning}
\ForEach{edge $e$ of $S$}{\label{lin:saving-use-s} \label{lin:saving-use-s}
	Let $S^{/e}$ be the graph after removing $e$ from $S$\label{lin:saving-use}\\
	Remove the values of the domains in $S$ that correspond to invalid assignments of $S^{/e}$\label{lin:saving-use-e}
} 

\BlankLine
\Comment{Unique labels}
%\If{$S$ has unique labels and $G$ has single labels}{ \label{lin:unique}
\If{$S$ and $G$ satisfy the unique labels optim. conditions}{ \label{lin:unique}
	\lIf{the size of any domain is less than $\tau$} {\Return $\mathit{false}$} \label{lin:unique1}
	\lElse \Return $\mathit{true}$ \label{lin:unique2}
}
\BlankLine
\Comment{Automorphisms}
Compute the automorphisms of $S$ \label{lin:auto-compute}

\BlankLine
\ForEach{variable $x$ and its domain $D$}{\label{lin:isFreq-var}
	$count \leftarrow 0$, $\mathit{timedoutSearch}\leftarrow\emptyset$\\
%	\BlankLine
%	\Comment{Automorphisms}
	\If{there is an automorphism with a computed domain $D'$} {\label{lin:auto-check}
		$D\leftarrow D'$ and move to the next $x$ variable (Line \ref{lin:isFreq-var})
	}
%	\BlankLine
	Apply arc consistency\\
	\lIf{the size of a domain is less than $\tau$} {\Return $\mathit{false}$}
	\BlankLine
	\Comment{Lazy search}
	\ForEach{element $u$ of $D$ }{\label{lin:isFreq-lazy}
		\lIf{$u$ is already marked}{$count$++}
		\Else{
			Search for a solution that assigns $u$ to $x$ for a given time threshold\label{lin:search1}\\
			\lIf{search timeouts}{\label{lin:timeout-start}
				Save the search state in a structure $\mathit{timedoutSearch}$
			}
			\If{a solution $\mathcal{S}ol$ is found}{
				Mark all values of $\mathcal{S}ol$ to the corresponding domains\\
				$count$++ 
			}
			\lElse{
				Remove $u$ from the domain $D$ and add $u$ to the invalid assignments of $D$ in $S$ \label{lin:saving-pop-1}
			}
			\lIf{$count=\tau$}{Move to the next variable (Line \ref{lin:isFreq-var}) \label{lin:isFreq-lazy-end}}
		}
	}

\BlankLine
\Comment{Resume timedout search if needed}
\If{$|\mathit{timedoutSearch}| + count \geq \tau$} { \label{lin:timedout-loop}
	\BlankLine
	\Comment{Decompose}
	Decompose graph $S$ into a set of graphs $\mathbf{Set}$ that contain the newly added edge \label{lin:decompose1}\\
	\lForEach{$s \in \mathbf{Set}$ }    {\label{lin:decompose2}
		Remove invalid assignments of $s$ from the respective domains of $S$
	}
      	\ForEach{$t \in \mathit{timedoutSearch}$ \label{lin:search-timeout}}    {
	Resume search from the saved state $t$\label{lin:search2}\\
	\If{a solution $\mathcal{S}ol$ is found}{
		Mark all values of $\mathcal{S}ol$ to the corresponding domains\\
		$count$++
	}
	\lElse{Remove $u$ from the domain $D$ and add $u$ to the invalid assignments of $D$ in $S$ \label{lin:saving-pop-2}}
	\lIf{$count=\tau$}{Move  to the next variable (Line \ref{lin:isFreq-var}) \label{lin:isFreq-var-2}}
      	}
}
  \Return $\mathit{false}$\ \ \Comment{Domain is exhausted and $count < \tau$}\label{lin:isFreq-end}
}
\Return $\mathit{true}$ 
\end{algorithm}

\commentg{Και οι άλλες σχέσεις, με αντίστοιχο τρόπο...}



\section{Εγκατάσταση και απαιτήσεις}
...

